% You can switch XeLaTeX/pdfLaTeX/LaTeX/LuaLaTeX in Settings
\documentclass[12pt, a4paper, oneside]{ctexart}
\usepackage{geometry} % 让一页显示的内容更多
\geometry{left=2.54cm, right=2.54cm, top=3.18cm, bottom=3.18cm}
\usepackage{amsmath, amsthm, amssymb} % 数学宏包
\usepackage{graphicx} % 放图片的宏包
% hyperref 要放在所有宏包的最后
\usepackage[bookmarks=true, colorlinks, citecolor=blue, linkcolor=black]{hyperref}

% 导言区

\title{我的第一个\LaTeX 文档}
\author{Mo}
\date{\today}

% 正文区

\begin{document}

\maketitle % 放置标题

\tableofcontents % 添加目录

\newpage
\section{一级标题}

\subsection{二级标题}

这里是正文. 如图~\ref{fig:kirigiri}所示 % ~是不可换行空格
，也如表~\ref{tab:my-table}所示。

\begin{figure}[htbp] % 自动选择最优位置
    \centering % 居中
    \includegraphics[width=0.3\textwidth]{图片.jpg}
    \caption{雾切响子} % 标题
    \label{fig:kirigiri} % 打标签，方便引用
\end{figure}

\begin{table}[htbp]
    \centering
    \caption{表格标题}
    \begin{tabular}{ccc}
        1 & 2 & 3 \\
        4 & 5 & 6 \\
        7 & 8 & 9
    \end{tabular}
    \label{tab:my-table}
\end{table}

\subsection{二级标题}

这里也是正文. 

%行内公式 inline
若$a>0$, $b>0$, 则$a+b>0$.
设$\lim_{n\to\infty}x_n=x$.

%行间公式 display
若$a>0$, $b>0$, 则
\[ % 也可用$$
a+b>0.
\]
设$\displaystyle\lim_{n\to\infty}x_n=x$.

\begin{enumerate}
    \item[(1)] 这是第一点; 
    \item[(2)] 这是第二点;
    \item 这是第三点. 
\end{enumerate}

\end{document}
